\section{集合的元素与子集}

\begin{enumerate}
\def\labelenumi{\arabic{enumi}.}
\item
  一个集合由一些元素组成，这些元素能够具有某些性质，并且能够彼此之间或与其它集合的元素之间具有某些关系。
\item
  在数学论证中，集合和元素用图形符号表示，这些符号通常是由字母（来自各种字母表）或字母与其它符号的组合构成。一个或多个集合的元素之间的关系，通过将这些表示元素的符号插入到该关系所特有的图式中来表示（†）；性质也是如此。
\end{enumerate}

一个字母可以表示一个集合的一个固定元素，也可以表示一个任意元素（也称为变量、自变量或通有元素）。当在一个关系（或性质）中，一个任意元素被一个固定元素替换时，称该任意元素被赋予这个固定元素作为其值。

为了指明一个未明确写出的关系中所出现的元素，我们用如下形式的记号来表示该关系： \(R\{x,y,z\}\) （如果 x,y,z 是所涉及的元素）。

\begin{enumerate}
\def\labelenumi{\arabic{enumi}.}
\setcounter{enumi}{2}
\item
  一个其中出现任意元素（*）的关系或性质，如果无论给这些任意元素赋予什么值，它都成为一个真命题，则称其为恒等式。如果 R 和 S 表示两个关系（或性质），那么当这些关系中所涉及的任意元素被固定为使得 R 为真的方式时，S 为真，则称 R 蕴含 S。关系（或性质）R 和 S 若彼此蕴含，则称它们等价。
\item
  让 \(R\{x, y, z\}\) 是变量 x, y, z 之间的关系。短语“对所有 x， \(R\{x, y, z\}\) ”是 y 和 z 之间的关系，如果对于 y 和 z 的这些给定值以及 x 的每一个值，R 都为真，则认为该关系对于 y 和 z 的给定值系统为真。短语“存在 x 使得 \(R\{x, y, z\}\) ”（或“对于某个 x， \(R\{x, y, z\}\) ”）同样是 y 和 z 之间的关系，如果在这些变量如此固定后，至少存在一个 x 的值使得 R 为真，则认为该关系对于 y 和 z 的给定值系统为真。对于任意多个变量之间的关系也是如此。
\end{enumerate}

如果 \(\overline{R}\) 表示 R 的否定，那么“对所有 x，R”的否定是“存在 x 使得 \(\overline{R}\) ”；“存在 x 使得 R”的否定是“对所有 x， \(\overline{R}\) ''.

\begin{enumerate}
\def\labelenumi{\arabic{enumi}.}
\setcounter{enumi}{4}
\item
  如果 R, S 表示两个关系，我们将“R 且 S”视为一个单一关系，当 R 和 S 都为真时，它被认为为真。同样，“R 或 S”是一个关系，当 R, S 中至少有一个为真时（特别是当它们都为真时），它被认为为真。因此，这里的“或”一词不具有它在日常用语中有时具有的排他性含义。设 \(\overline{R}\) , \(\overline{S}\) 分别表示 R, S 的否定。那么“R 且 S”的否定是“ \(\overline{R}\) 要么 \(\overline{S}\) ”，而“R 或 S”的否定是“ \(\overline{R}\) 并且 \(\overline{S}\) ''.
\item
  通过在符号“=”（读作“等于”）的两侧各写一个符号，我们得到一个称为相等关系的关系，它表示这两个符号代表同一个元素。该关系的否定是通过在符号“≠”（读作“不等于”或“不同于”）的两侧写相同的符号而得到的。
\item
  给定一个集合 E 和 E 的一个通有元素的一个性质，E 中具有该性质的那些元素构成一个新集合，称为 E 的一个子集。因此，两个等价的性质定义 E 的同一个子集，反之亦然。
\end{enumerate}

设 A 是 E 的一个子集。当 x 是 E 的一个通有元素时，性质“x 属于 A”（即“x 是 A 的一个元素”）写作“ \(x \in A\) ”；具有该性质的元素的集合显然就是 A。

该性质的否定写作“\(x \notin A\)”（读作“x 不属于 A”）；E 中具有该性质的元素的集合称为 A 的补集，记作 CA 或 E --- A。

\begin{enumerate}
\def\labelenumi{\arabic{enumi}.}
\setcounter{enumi}{7}
\tightlist
\item
  某些性质，例如 \(x = x\) 对E的所有元素都为真。任意两个这样的性质是等价的，并且它们所定义的子集就是集合E本身。
\end{enumerate}

另一方面，某些性质，例如 \(x \neq x\) 对E的任何元素都不为真。同样，任意两个这样的性质是等价的，并且它们所定义的子集称为E的空子集，记为 \(\phi\) .

注意，E和 \(\phi\) 互为补集。

\begin{enumerate}
\def\labelenumi{\arabic{enumi}.}
\setcounter{enumi}{8}
\item
  设 \(a\) 是 \(E\) 的一个确定元素。某些性质，例如 \(x = a\) ，仅对单个元素 \(a\) 为真。任意两个这样的性质是等价的；它们所定义的子集记为 \(\{a\}\) ，并称为由 \(a\) 单独组成的子集。
\item
  以集合E的所有子集为元素的集合称为E的子集族，记为 \(\mathfrak{P}(\mathbf{E})\) 。我们有 \(\phi \in \mathfrak{P}(\mathbf{E})\) , \(\mathbf{E} \in \mathfrak{P}(\mathbf{E})\) ，以及 \(\{x\} \in \mathfrak{P}(\mathbf{E})\) 对于所有 \(x \in E\) . 如果x表示E的一个一般元素，X表示 \(\mathfrak{P}(\mathbf{E})\) 的一个一般元素，则x与X之间的关系“ \(x \in X\) ”称为属于关系。
\item
  设 \(x\) 和 \(y\) 为 \(\mathbf{E}\) 的两个元素，并令 \(\mathbf{X}\) 是 \(\mathfrak{P}(\mathbf{E})\) 的一个一般元素。那么等式“ \(x = y\) ”等价于关系“对所有 \(\mathbf{X}\) 使得 \(x \in \mathbf{X}\) ，我们有 \(y \in \mathbf{X}\) ''.
\item
  设X，Y是集合E的两个子集。如果性质 \(x \in X\) 蕴含性质 \(x \in Y\) ，换句话说，如果X的每个元素都属于Y，那么我们说X包含于Y，或者说Y包含X，或者说X是Y的子集。X与Y之间的这种关系称为包含关系（X包含于Y），记为“X ⊂ Y”或“Y ⊃ X”。其否定记为“X ⊄ Y”或“Y ⊅ X”。
\end{enumerate}

对于E的所有子集X，我们有 \(\phi \subset X\) 并且 \(X \subset E\) 。属于关系“ \(x \in X\) ”等价于“ \(\{x\} \subset X\) ''.

关系“X ⊂ Y且Y ⊂ Z”蕴含“X ⊂ Z”。

关系“X ⊂ Y”不排除“X = Y”的可能性。关系“X ⊂ Y且Y ⊂ X”等价于“X = Y”。

\begin{enumerate}
\def\labelenumi{\arabic{enumi}.}
\setcounter{enumi}{12}
\tightlist
\item
  设X和Y是E的任意两个子集。所有具有性质“ \(x \in X\) 要么 \(x \in Y\) ”的元素所组成的集合记为X ∪ Y，称为X与Y的并集。所有具有性质“ \(x \in X\) 并且 \(x \in Y\) ”的元素所组成的集合记为X ∩ Y，称为X与Y的交集。
\end{enumerate}

若干E的子集的并集和交集以相同方式定义。

如果 \(x, y, z\) 是 \(\mathbf{E}\) 的三个元素，并集 \(\{x\} \cup \{y\} \cup \{z\}\) 记作 \(\{x, y, z\}\) 。类似地，对于任意数量（各自命名的）元素。

设X和Y是E的两个子集。根据 \(X \cap Y \neq \phi\) 要么 \(X \cap Y = \phi\) ，我们称X与Y相交或不相交。

\begin{enumerate}
\def\labelenumi{\arabic{enumi}.}
\setcounter{enumi}{13}
\tightlist
\item
  在以下命题的陈述中，X、Y、Z表示同一集合E的任意子集。
\end{enumerate}

\begin{enumerate}
\def\labelenumi{(\alph{enumi})}
\item
  我们有 \(\phi = \mathbb{C}\mathbf{E},\quad \mathbf{E} = \mathbb{C}\phi .\)
\item
  对于所有 \(\mathbf{X}\) 我们拥有
\end{enumerate}

\begin{enumerate}
\def\labelenumi{(\arabic{enumi})}
\tightlist
\item
\end{enumerate}

\[
\mathbb {C} (\mathbb {C} \mathbf {X}) = \mathbf {X};\tag{2}
\]

\[
\mathbf {X} \cup \mathbf {X} = \mathbf {X}, \quad \mathbf {X} \cap \mathbf {X} = \mathbf {X};\tag{3}
\]

\[
\mathbf {X} \cup (\mathbb {C X}) = \mathbf {E}, \quad \mathbf {X} \cap (\mathbb {C X}) = \varnothing ;\tag{4}
\]

\[
\mathbf {X} \cup \phi = \mathbf {X}, \quad \mathbf {X} \cap \mathbf {E} = \mathbf {X};\tag{5}
\]

\[
\mathbf {X} \cup \mathbf {E} = \mathbf {E}, \quad \mathbf {X} \cap \phi = \phi .
\]

\begin{enumerate}
\def\labelenumi{(\alph{enumi})}
\setcounter{enumi}{2}
\tightlist
\item
  对于所有 \(\mathbf{X},\mathbf{Y}\) 我们拥有
\end{enumerate}

\begin{enumerate}
\def\labelenumi{(\arabic{enumi})}
\setcounter{enumi}{5}
\tightlist
\item
\end{enumerate}

\[
\mathbf {X} \cup \mathbf {Y} = \mathbf {Y} \cup \mathbf {X}, \quad \mathbf {X} \cap \mathbf {Y} = \mathbf {Y} \cap \mathbf {X} \quad (\text { 交换律 });\tag{7}
\]

\[
\mathbf {X} \subset \mathbf {X} \cup \mathbf {Y}, \quad \mathbf {X} \cap \mathbf {Y} \subset \mathbf {X};\tag{8}
\]

\[
\mathbf {C} (\mathbf {X} \cup \mathbf {Y}) = (\mathbf {C X}) \cap (\mathbf {C Y}), \quad \mathbf {C} (\mathbf {X} \cap \mathbf {Y}) = (\mathbf {C X}) \cup (\mathbf {C Y}).
\]

\begin{enumerate}
\def\labelenumi{(\alph{enumi})}
\setcounter{enumi}{3}
\item
  关系 \(\mathbf{X} \subset \mathbf{Y}\) , \(\complement \mathbf{X} \supseteq \complement \mathbf{Y}\) , \(\mathbf{X} \cup \mathbf{Y} = \mathbf{Y}\) , \(\mathbf{X} \cap \mathbf{Y} = \mathbf{X}\) 是等价的。
\item
  关系 \(\mathbf{X} \cap \mathbf{Y} = \varnothing\) , \(\mathbf{X} \subset \mathbb{C}\mathbf{Y}\) , \(\mathbf{Y} \subset \mathbb{C}\mathbf{X}\) 是等价的。
\item
  关系 \(\mathbf{X} \cup \mathbf{Y} = \mathbf{E}\) , \(\complement \mathbf{X} \subset \mathbf{Y}\) , \(\complement \mathbf{Y} \subset \mathbf{X}\) 是等价的。
\item
  对于所有X、Y、Z，我们有
\end{enumerate}

\[
\mathbf {X} \cup (\mathbf {Y} \cup \mathbf {Z}) = (\mathbf {X} \cup \mathbf {Y}) \cup \mathbf {Z} = \mathbf {X} \cup \mathbf {Y} \cup \mathbf {Z}\tag{9}
\]

\[
\mathrm{X} \cap (\mathrm{Y} \cap \mathrm{Z}) = (\mathrm{X} \cap \mathrm{Y}) \cap \mathrm{Z} = \mathrm{X} \cap \mathrm{Y} \cap \mathrm{Z}\tag{10}
\]

\[
\mathrm{X} \cup (\mathrm{Y} \cap \mathrm{Z}) = (\mathrm{X} \cup \mathrm{Y}) \cap (\mathrm{X} \cup \mathrm{Z})
\]

\[
\mathbf {X} \cap (\mathbf {Y} \cup \mathbf {Z}) = (\mathbf {X} \cap \mathbf {Y}) \cup (\mathbf {X} \cap \mathbf {Z}) \Big \}
\]

（分配律）。

\begin{enumerate}
\def\labelenumi{(\alph{enumi})}
\setcounter{enumi}{7}
\item
  关系“X ⊂ Y”蕴含关系“X ∪ Z ⊂ Y ∪ Z”和“X ∩ Z ⊂ Y ∩ Z”。
\item
  关系“Z ⊂ X且Z ⊂ Y”等价于“Z ⊂ X ∩ Y”。关系“X ⊂ Z且Y ⊂ Z”等价于“X ∪ Y ⊂ Z”。
\end{enumerate}

\begin{enumerate}
\def\labelenumi{\arabic{enumi}.}
\setcounter{enumi}{14}
\tightlist
\item
  由恒等式(8)我们得出结论：如果E的子集A是通过仅对E的其他子集X、Y、Z应用运算而得到的， \(\mathbb{C}\) , \(\cup\) , \(\cap\) （以任意顺序），则补集 \(\mathbb{C}\mathbf{A}\) 可通过将子集X、Y、Z分别替换为其补集，并相应调整运算来得到。 \(\cup\) , \(\cap\) 由 \(\cap\) , \(\cup\) 分别地，同时保持运算的顺序。这就是对偶规则。给定上述形式子集之间的等式 A = B，考虑等价等式 \(\mathbb{C}\mathbf{A} = \mathbb{C}\mathbf{B}\) 如果我们用对偶规则得到的表达式替换 \(\mathbb{C}\mathbf{A}\) 并且 \(\mathbb{C}\mathbf{B}\) ，然后再将 \(\mathbb{C}\mathbf{X}\) , \(\mathbb{C}\mathbf{Y}\) , \(\mathbb{C}\mathbf{Z}\) 分别替换为 X、Y、Z，反之亦然，我们就得到一个称为 A = B 的对偶的等式。对于包含关系 A ⊂ B，我们也可以做同样的操作，但此时必须注意将符号“⊂”替换为“⊃”。
\end{enumerate}

上面编号相同的恒等式互为对偶。

\begin{enumerate}
\def\labelenumi{\arabic{enumi}.}
\setcounter{enumi}{15}
\tightlist
\item
  在某些问题中，我们必须考虑集合E的一个固定子集A。如果X是E的任意子集，则集合 \(A \cap X\) 称为X在A上的迹，有时记为 \(X_{A}\) ；在这种情况下，它被视为A的一个子集。对于E的所有子集X，Y，我们有
\end{enumerate}

\[
\begin{array}{c} (\mathrm{X} \cup \mathrm{Y}) _ {\mathrm{A}} = \mathrm{X} _ {\mathrm{A}} \cup \mathrm{Y} _ {\mathrm{A}}, \quad (\mathrm{X} \cap \mathrm{Y}) _ {\mathrm{A}} = \mathrm{X} _ {\mathrm{A}} \cap \mathrm{Y} _ {\mathrm{A}}, \\ \mathbb {C} _ {\mathrm{A}} \mathrm{X} _ {\mathrm{A}} = (\mathbb {C} _ {\mathrm{E}} \mathrm{X}) _ {\mathrm{A}}, \end{array}
\]

其中 \(C_{E}X\) 表示X在E中的补集，并且 \(C_{A}X_{A}\) 表示 \(X_{A}\) 在 A 中的补集。

如果 \(\mathcal{E}\) 表示 \(\mathbf{E}\) 的子集族，则集合 \(\mathcal{E}_{\mathbf{A}}\) （ \(\mathbf{A}\) 与集合族 \(\mathcal{E}\) 中各集合的交的全体）称为 \(\mathcal{E}\) 在 \(\mathbf{A}\) 上的迹.

